\documentclass[11pt]{article}
\usepackage[margin=1in]{geometry}
\usepackage{amsmath,amssymb,amsfonts,mathtools}
\usepackage{array,booktabs,enumitem,graphicx,xcolor,hyperref}
\usepackage[T1]{fontenc}
\usepackage[utf8]{inputenc}
\title{Multiple Crossover Phenomena and Scale Hopping in Two Dimensions}
\author{Source-backed arXiv sample}
\date{}
\begin{document}
\maketitle
\section{Extracted Lists}

The list environments below are extracted from arXiv LaTeX source and wrapped for pdfLaTeX validation.

\begin{enumerate}

\item Positions of the nearly marginal thermal operators in the Kac table of a
unitary  minimal model $M_p$ for $p \gg 1$. The operators
$\phi_{(n,n+2)}$ are  relevant, while the operators
$\phi_{(n,n-2)}$ and $\tilde \phi_{(n,n)}$ are irrelevant.

\item Special solutions of the RG equations in the vicinity of $M_{p,p-1}$
(schematic).
(a) A trajectory in ${\cal R}_p$.
(b) A trajectory in ${\cal C}_{p-1}$.
(c) A redundant trajectory.

\item Self-conjugate trajectories in the vicinity of the fixed point $M_p$
(schematic).
(a) A trajectory in $M_{p+1,p-1}$.
(b) A self-similar trajectory.

\item The $C$-function $C(\tau,\tau_0)$ of the unique self-similar trajectory
$\check u_{p}^i(\tau,s(\tau_0))$: a step in the staircase pattern for   $\tau_0
= 3.2$, $3.6$, $4.0$ and $4.4$ (solid lines). For larger values of  $\tau_0$,
the steps get more pronounced as the solutions tend towards the limit
trajectory $M_{p,p-1}$ (long-dashed line).

\item
(a) The RG flow in the vicinity of $M_p$. A self-similar trajectory of regime 1
(solid line) visits all fixed points, trajectories in regime 2 or 3
(long-dashed lines) visit two fixed points, and a trajectory in regime 4
(short-dashed line) visits only one fixed point.
(b) The resulting phase diagram in the $(t_{p}, \bar t_{p})$-plane.

\end{enumerate}
\end{document}
